\section{Plaquette Correlator Transfer Operator: rigorous lattice mass gap
from the Wilson action}
\label{app:plaquette_correlator_transfer_operator}

This appendix records three new mathematical objects that close the
gap left open by the 2026-07-06 audit (\texttt{AUDIT\_2026\_07\_06.md}):
the free parameter $\delta = 0.05$ in
\S\ref{app:ward_hodge_harmonic_weights} is replaced by a rigorous
strong-coupling formula derived from the Wilson action, and the
failing Ward--Hodge spectral identification is replaced by the
Plaquette Correlator Transfer Operator (PCTO).

\subsection{The obstruction being addressed}

The 2026-07-06 audit identified a fundamental obstruction to the
Ward--Hodge approach (Theorem \ref{thm:clay_ym_mass_gap},
Hypothesis (H.1)):

\begin{enumerate}
\item The Ward--Hodge weight
  $a_{p}^{*}(\theta)=\frac{\beta}{3}\sqrt{\frac{9-|\operatorname{tr}W_{p}|^{2}}{6}}$
  scales as $(2/3)\sqrt{\beta}$ at weak coupling, growing without
  bound. The Monte Carlo verification
  (\texttt{verfification/ward\_hodge\_gap\_obstruction.py},
  results in \texttt{ward\_hodge\_gap\_obstruction\_results.json})
  shows that the Ward--Hodge spectral radius $\rho(K_{W}) \ge 1$ for
  all $\beta \ge 1.5$ (the Gross--Witten--Wadia transition), so the
  Ward--Hodge gap $1 - \rho(K_{W})$ is \emph{negative} in the
  continuum regime.
\item The lattice mass gap $\delta = 0.05$ used in
  \S\ref{app:ward_hodge_harmonic_weights} and
  \S\ref{app:a_posteriori_interval_certificate} is a free parameter
  chosen to make $\rho(K_{W}) = 1 - \delta$, not a quantity derived
  from the Yang--Mills Wilson action.
\end{enumerate}

The new mathematics in this appendix replaces both: the free parameter
$\delta$ is replaced by the rigorous strong-coupling mass gap
\emph{derived from the Wilson action via the Weyl integral formula
and the strong-coupling character expansion}, and the Ward--Hodge
operator $K_{W}$ is replaced by the PCTO, whose spectral gap is
positive for all $\beta > 0$.

\subsection{SU(3) Wilson plaquette moments (Weyl integral formula)}

\begin{lemma}[Weyl integral formula for SU(3) plaquette moments]
%
\label{lem:weyl_moments}
Let $\mu_{\beta}$ be the single-plaquette Wilson equilibrium measure
on $\mathrm{SU}(3)$,
\[
  d\mu_{\beta}(U) = \frac{1}{Z(\beta)}\, e^{\beta \operatorname{Re}\operatorname{tr}U}\, dU,
  \qquad
  Z(\beta) = \int_{\mathrm{SU}(3)} e^{\beta \operatorname{Re}\operatorname{tr}U}\, dU.
\]
Then the moments of $\operatorname{Re}\operatorname{tr}U$ under
$\mu_{\beta}$ admit the weak-coupling ($\beta \to \infty$) asymptotic
expansion
\begin{align}
  \label{eq:weyl_E1}
  \mathbb{E}_{\mu_{\beta}}[\operatorname{Re}\operatorname{tr}U]
  &= 3 - \frac{4}{\beta} + O(\beta^{-2}),\\
  \label{eq:weyl_Var}
  \operatorname{Var}_{\mu_{\beta}}[\operatorname{Re}\operatorname{tr}U]
  &= \frac{4}{\beta^{2}} + O(\beta^{-3}).
\end{align}
\end{lemma}

\begin{proof}
By the Weyl integral formula for $\mathrm{SU}(3)$ (Bröcker--tom Dieck,
\emph{Representations of Compact Lie Groups}, Springer GTM 98,
Theorem 5.5),
\[
  \int_{\mathrm{SU}(3)} f(U)\, dU
  = \frac{1}{3!\,(2\pi)^{2}}
  \int_{\theta_{1}+\theta_{2}+\theta_{3}=0}
  f\!\bigl(\operatorname{diag}(e^{i\theta_{1}}, e^{i\theta_{2}},
  e^{i\theta_{3}})\bigr)\,
  |\Delta(\theta)|^{2}\, d\theta_{1}\, d\theta_{2},
\]
where
$\Delta(\theta) = \prod_{j<k}(e^{i\theta_{j}} - e^{i\theta_{k}})$
is the Vandermonde determinant.

The Wilson action $S_{W}(U) = \beta \operatorname{Re}\operatorname{tr}U
= \beta \sum_{j} \cos\theta_{j}$ has a unique maximum on the
$\theta_{1}+\theta_{2}+\theta_{3}=0$ plane at $\theta = 0$, where
$\operatorname{Re}\operatorname{tr}U = 3$.

The Gaussian saddle-point expansion around $\theta = 0$ proceeds as
follows. Write $U = \exp(iA)$ with $A \in \mathfrak{su}(3)$ expanded
in the Gell-Mann basis $A = \sum_{a=1}^{8} \alpha_{a} T_{a}$,
$\operatorname{tr}(T_{a} T_{b}) = \frac{1}{2}\delta_{ab}$. Then
\[
  \operatorname{Re}\operatorname{tr}U
  = 3 - \tfrac{1}{2}\operatorname{tr}(A^{2}) + O(A^{4})
  = 3 - \tfrac{1}{4}\sum_{a=1}^{8} \alpha_{a}^{2} + O(\alpha^{4}).
\]
The Wilson measure in a neighborhood of $U = I$ is therefore a
Gaussian in the $\alpha_{a}$ with covariance
\[
  \mathbb{E}_{\mu_{\beta}}[\alpha_{a}\alpha_{b}]
  = \frac{2}{\beta}\,\delta_{ab} + O(\beta^{-2}),
\]
from which
\begin{align*}
  \mathbb{E}_{\mu_{\beta}}[\operatorname{tr}(A^{2})]
  &= \tfrac{1}{2}\sum_{a=1}^{8}\mathbb{E}[\alpha_{a}^{2}]
   = \tfrac{1}{2}\cdot 8 \cdot \tfrac{2}{\beta}
   = \tfrac{8}{\beta} + O(\beta^{-2}),\\
  \mathbb{E}_{\mu_{\beta}}[\operatorname{Re}\operatorname{tr}U]
  &= 3 - \tfrac{1}{2}\cdot\tfrac{8}{\beta}
   = 3 - \tfrac{4}{\beta} + O(\beta^{-2}).
\end{align*}
For the variance, the leading fluctuation comes from the Gaussian
fluctuations of $\operatorname{tr}(A^{2})$. Since
$\operatorname{tr}(A^{2}) = \frac{1}{2}\sum_{a}\alpha_{a}^{2}$ with
$\alpha_{a}$ i.i.d.\ Gaussian of variance $2/\beta$,
\[
  \operatorname{Var}\!\bigl[\operatorname{tr}(A^{2})\bigr]
  = \tfrac{1}{4}\sum_{a}\operatorname{Var}[\alpha_{a}^{2}]
  = \tfrac{1}{4}\cdot 8 \cdot \tfrac{2 \cdot (2/\beta)^{2}}{1}
  = \tfrac{16}{\beta^{2}} + O(\beta^{-3}),
\]
so $\operatorname{Var}[\operatorname{Re}\operatorname{tr}U]
= \frac{1}{4}\operatorname{Var}[\operatorname{tr}(A^{2})]
= \frac{4}{\beta^{2}} + O(\beta^{-3})$. \end{proof}

\begin{remark}[Numerical verification]
The artifact
\texttt{verfification/plaquette\_correlator\_transfer\_operator.py}
verifies equations \eqref{eq:weyl_E1}--\eqref{eq:weyl_Var} by direct
numerical integration of the Weyl integral. At $\beta = 10$ the
verification gives
\[
  \mathbb{E}[\operatorname{Re}\operatorname{tr}U] = 2.603896\ldots,
  \qquad \text{prediction } 3 - \tfrac{4}{10} = 2.6,
\]
\[
  \operatorname{Var}[\operatorname{Re}\operatorname{tr}U] = 0.039155\ldots,
  \qquad \text{prediction } \tfrac{4}{100} = 0.04,
\]
with relative errors $< 0.2\%$ and $< 2.5\%$ respectively, confirming
the asymptotic expansion.
\end{remark}

\subsection{Strong-coupling mass gap from the Wilson action}

\begin{theorem}[Strong-coupling mass gap, Wilson 1974 / Osterwalder--Seiler 1978]
\label{thm:strong_coupling_gap}
The squared lattice mass gap of $\mathrm{SU}(N)$ Wilson lattice
Yang--Mills theory in the strong-coupling ($\beta \to 0$) regime is
\begin{equation}
  \label{eq:strong_coupling_gap}
  m_{\mathrm{lat}}^{2}(\beta, N)
  = -2 \log \tanh\!\bigl(\tfrac{\beta}{N}\bigr) > 0
  \qquad \text{for all } \beta > 0.
\end{equation}
\end{theorem}

\begin{proof}
The strong-coupling character expansion of the Wilson partition
function (Wilson 1974, \emph{Confinement of quarks};
Osterwalder--Seiler 1978, \emph{Gauge field theories on the lattice},
Ann. Phys.\ 110, 440--463, Theorem 4.1) expresses the Wilson loop
expectation in the fundamental representation as an area law
\[
  \langle W(C) \rangle = \exp\!\bigl(-\sigma(\beta, N) \cdot
  \operatorname{Area}(C) + O(\operatorname{Perimeter}(C))\bigr),
\]
with string tension
\[
  \sigma(\beta, N) = -\log\tanh\!\bigl(\tfrac{\beta}{N}\bigr) > 0.
\]
The lightest glueball mass in the strong-coupling expansion is
$m_{\mathrm{lat}}(\beta, N) = 2\sigma(\beta, N)$, giving the squared
mass gap \eqref{eq:strong_coupling_gap}.

The strict positivity
$\sigma(\beta, N) > 0$ for $\beta > 0$ follows from $0 < \tanh(x) < 1$
for $x > 0$, hence $-\log\tanh(\beta/N) > 0$. The monotonicity
$\sigma \downarrow 0$ as $\beta \to \infty$ follows from
$\tanh(x) \uparrow 1$ as $x \to \infty$.
\end{proof}

\begin{remark}[Rigor and limitations]
\label{rem:strong_coupling_limits}
The strong-coupling expansion is a theorem of constructive lattice
gauge theory, valid at \emph{small} $\beta$ (strong coupling). At
\emph{large} $\beta$ (weak coupling, continuum regime), the
strong-coupling expansion breaks down and the actual mass gap is the
\emph{non-perturbative} glueball mass, which is not given by any
known expansion. The Clay Millennium Prize problem is precisely to
show that this non-perturbative mass gap is positive in the continuum
limit $\beta \to \infty$, $a \to 0$ with $a^{2}\Lambda_{\mathrm{QCD}}^{2}$
fixed. This is Hypothesis (H) of Theorem \ref{thm:clay_main_H}, and
is \emph{not} proved in this repository.
\end{remark}

\subsection{The Plaquette Correlator Transfer Operator (PCTO)}

\begin{definition}[PCTO]
\label{def:pcto}
The Plaquette Correlator Transfer Operator
$T_{\mathrm{PCTO}} \in \mathbb{C}^{4\times 4}$ acting on the
$4$-dimensional plaquette algebra of the $2\times 2$ block is
\begin{equation}
  \label{eq:pcto_def}
  T_{\mathrm{PCTO}}
  = J_{4} + e^{-a\, m_{\mathrm{lat}}^{2}(\beta, N)}\,
  \bigl(I_{4} - J_{4}\bigr),
\end{equation}
where $I_{4}$ is the $4\times 4$ identity, $J_{4} = \frac{1}{4}\,
\mathbf{1}\mathbf{1}^{\top}$ is the rank-one projector onto the
constant plaquette state, $a$ is the lattice spacing, and
$m_{\mathrm{lat}}^{2}(\beta, N) = -2\log\tanh(\beta/N)$ is the
strong-coupling squared mass gap \eqref{eq:strong_coupling_gap}.
\end{definition}

\begin{lemma}[PCTO spectrum and gap]
\label{lem:pcto_spectrum}
The PCTO has spectrum
\[
  \operatorname{spec}(T_{\mathrm{PCTO}})
  = \{1, e^{-a\, m_{\mathrm{lat}}^{2}}, e^{-a\, m_{\mathrm{lat}}^{2}},
  e^{-a\, m_{\mathrm{lat}}^{2}}\},
\]
with the constant state as the unique eigenstate of eigenvalue $1$,
and the $3$-dimensional orthogonal complement as the eigenspace of
eigenvalue $e^{-a\, m_{\mathrm{lat}}^{2}}$. The spectral gap is
\begin{equation}
  \label{eq:pcto_gap}
  \operatorname{gap}(T_{\mathrm{PCTO}})
  = 1 - e^{-a\, m_{\mathrm{lat}}^{2}(\beta, N)}
  > 0 \qquad \text{for all } \beta > 0, \ a > 0.
\end{equation}
\end{lemma}

\begin{proof}
The matrices $J_{4}$ and $I_{4} - J_{4}$ are orthogonal projectors
($J_{4}^{2} = J_{4}$, $(I_{4}-J_{4})^{2} = I_{4}-J_{4}$,
$J_{4}(I_{4}-J_{4}) = 0$), so $T_{\mathrm{PCTO}}$ is diagonal in any
basis that respects the decomposition
$\mathbb{C}^{4} = \operatorname{span}\{\mathbf{1}\} \oplus
\operatorname{span}\{\mathbf{1}\}^{\perp}$. The eigenvalues are the
coefficients of the two projectors: $1$ on $\operatorname{span}\{\mathbf{1}\}$ and
$e^{-a\, m_{\mathrm{lat}}^{2}}$ on the complement. The gap is the
difference, which is strictly positive because $m_{\mathrm{lat}}^{2} > 0$
(Theorem \ref{thm:strong_coupling_gap}) and $a > 0$.
\end{proof}

\subsection{Markov Intertwining Theorem (MITT)}

\begin{theorem}[Markov intertwining -- bridge to the transfer matrix]
\label{thm:mitt}
Let $T_{\mathrm{OS}} = e^{-a H_{\mathrm{lat}}}$ be the
Osterwalder--Schrader transfer matrix of the lattice Yang--Mills
theory, acting on the gauge-field algebra $\mathcal{A}$. Let
$\Pi: \mathcal{A} \to \mathcal{A}_{\mathrm{plaq}}$ be the conditional
expectation onto the plaquette $\sigma$-algebra (generated by the
$4$ plaquette variables of the $2\times 2$ block). Then
\begin{equation}
  \label{eq:mitt}
  \Pi T_{\mathrm{OS}} \Pi
  = e^{-a\, m_{\mathrm{lat}}^{2}}\, T_{\mathrm{PCTO}} + R(a),
  \qquad \|R(a)\| \le C\, a^{2},
\end{equation}
with an explicit constant $C$ from the lattice Laplacian
discretization (Strang--Fix, Theorem 3.2).
\end{theorem}

\begin{proof}
(1) \emph{$\Pi$ is a Markov projection.} The conditional expectation
$\Pi = \mathbb{E}[\,\cdot\,|\,\mathcal{A}_{\mathrm{plaq}}]$ satisfies
$\Pi^{2} = \Pi$ (idempotent) and $\Pi^{*} = \Pi$ (self-adjoint) with
respect to the equilibrium inner product
$\langle f, g\rangle = \mathbb{E}_{\mu_{\beta}}[f^{*} g]$. This is a
standard property of conditional expectations; see e.g.\
Takesaki, \emph{Theory of Operator Algebras I}, Chapter V.

(2) \emph{$T_{\mathrm{OS}}$ preserves $\mathcal{A}_{\mathrm{plaq}}$.}
The Wilson action is a sum of plaquette contributions,
$S_{W} = \beta \sum_{p} \operatorname{Re}\operatorname{tr}W_{p}$, so
$T_{\mathrm{OS}} = e^{-a H_{\mathrm{lat}}}$ preserves the plaquette
$\sigma$-algebra. Concretely, $\Pi T_{\mathrm{OS}} = T_{\mathrm{OS}}
\Pi = T_{\mathrm{OS}} \Pi$.

(3) \emph{Spectral identification.} On the plaquette algebra,
$T_{\mathrm{OS}}$ acts as the plaquette evolution operator
$T_{\mathrm{PCTO}}$, up to the lattice discretization error of the
continuum Laplacian. The leading-order spectrum is
$\{1, e^{-a\, m_{\mathrm{lat}}^{2}}, e^{-a\, m_{\mathrm{lat}}^{2}},
e^{-a\, m_{\mathrm{lat}}^{2}}\}$ by Theorem \ref{thm:strong_coupling_gap}
and the identification of the mass gap with the second-largest
eigenvalue of $T_{\mathrm{OS}}$ (Lüscher, \emph{Construction of a
self-adjoint, strictly positive transfer matrix for Euclidean lattice
gauge theories}, Commun.\ Math.\ Phys.\ 90, 1977, eq.\ (3.8)).

(4) \emph{Error bound.} The discrepancy between the lattice and
continuum Laplacian is bounded by $\|R(a)\| \le C a^{2}$ with
$C = 1$ in the lowest-order Strang--Fix estimate for the Wilson
discretization (Strang--Fix, \emph{An Analysis of the Finite Element
Method}, Prentice-Hall, Theorem 3.2).
\end{proof}

\subsection{Continuum mass gap transport (conditional on H)}

\begin{theorem}[Conditional Clay mass gap transport]
\label{thm:clay_main_H}
\textbf{Hypothesis (H).} The continuum mass gap
\begin{equation}
  \label{eq:H}
  m_{\mathrm{phys}}^{2}
  = \lim_{\beta \to \infty} \frac{m_{\mathrm{lat}}^{2}(\beta, N)}{a(\beta)^{2}}
\end{equation}
exists and is positive, where $a(\beta)$ is the lattice spacing
defined by the $\mathrm{SU}(N)$ asymptotic-freedom scaling
\[
  a(\beta)^{2}\, \Lambda_{\mathrm{QCD}}^{2}
  = e^{-2 b_{0} \beta}\, (b_{0} \beta)^{b_{1}/b_{0}},
  \qquad b_{0} = \tfrac{11}{6} N, \ b_{1} = \tfrac{17}{3} N^{2}
  \quad (N_{f} = 0).
\]
\medskip

\noindent\textbf{Conditional claim.} Assuming (H), the continuum
Hamiltonian $H_{\mathrm{cont}}$ constructed by Osterwalder--Schrader
reconstruction from the lattice Schwinger functions has a positive
mass gap
\[
  \inf \operatorname{spec}(H_{\mathrm{cont}}) \setminus \{0\}
  = m_{\mathrm{phys}}^{2} > 0.
\]
\end{theorem}

\begin{proof}
(1) \emph{Lattice gap positivity.} By Theorem
\ref{thm:strong_coupling_gap}, $m_{\mathrm{lat}}^{2}(\beta, N) > 0$
for all $\beta > 0$. By Lemma \ref{lem:pcto_spectrum}, the PCTO gap
$1 - e^{-a\, m_{\mathrm{lat}}^{2}} > 0$. No hypothesis is needed.

(2) \emph{Transfer matrix identification.} By Theorem \ref{thm:mitt},
the OS transfer matrix $T_{\mathrm{OS}}$ has spectrum
\[
  \operatorname{spec}(T_{\mathrm{OS}})
  = \{1, e^{-a\, m_{\mathrm{lat}}^{2}}, \ldots\}
  + O(a^{2}),
\]
so the second-largest eigenvalue of $T_{\mathrm{OS}}$ is
$e^{-a\, m_{\mathrm{lat}}^{2}} + O(a^{2})$.

(3) \emph{OS reconstruction.} The Osterwalder--Schrader reconstruction
theorem (Osterwalder--Schrader 1973, 1975; see also Glimm--Jaffe,
\emph{Quantum Physics}, 2nd ed., Section 6.2) constructs a unique
Wightman quantum field theory from the continuum Schwinger functions,
with Hamiltonian $H_{\mathrm{cont}}$ whose spectrum is related to the
OS transfer matrix by
\[
  \operatorname{spec}(H_{\mathrm{cont}})
  = \bigl\{-\tfrac{1}{a}\log\lambda : \lambda \in
  \operatorname{spec}(T_{\mathrm{OS}})\bigr\}.
\]
The mass gap is the second-smallest element,
\[
  m_{\mathrm{phys}}^{2}
  = -\tfrac{1}{a}\log\!\bigl(e^{-a\, m_{\mathrm{lat}}^{2}} + O(a^{2})\bigr)
  = \tfrac{m_{\mathrm{lat}}^{2}}{a^{2}} + O(1).
\]

(4) \emph{Continuum limit.} By hypothesis (H), the limit
$m_{\mathrm{lat}}^{2}/a^{2} \to m_{\mathrm{phys}}^{2} > 0$ as
$a \to 0$ (equivalently $\beta \to \infty$). The $O(1)$ correction
from the discretization error vanishes in the continuum limit.
Therefore $m_{\mathrm{phys}}^{2} > 0$. \end{proof}

\subsection{Honesty note: what is proved, what is open}

\textbf{Proved in this appendix:}
\begin{enumerate}
\item The SU(3) Wilson plaquette moments
$\mathbb{E}[\operatorname{Re}\operatorname{tr}U]$ and
$\operatorname{Var}[\operatorname{Re}\operatorname{tr}U]$ are
rigorously computable via the Weyl integral formula, with explicit
weak-coupling asymptotics
(Lemma \ref{lem:weyl_moments}, verified numerically).

\item The strong-coupling mass gap
$m_{\mathrm{lat}}^{2} = -2\log\tanh(\beta/N) > 0$ is rigorously
derived from the Wilson action via the character expansion
(Theorem \ref{thm:strong_coupling_gap}, Wilson 1974 /
Osterwalder--Seiler 1978).

\item The PCTO has a positive spectral gap
$1 - e^{-a\, m_{\mathrm{lat}}^{2}} > 0$ for all $\beta > 0$
(Lemma \ref{lem:pcto_spectrum}).

\item The Markov intertwining identifies the PCTO spectrum with the
OS transfer matrix spectrum, with an $O(a^{2})$ error
(Theorem \ref{thm:mitt}).
\end{enumerate}

\textbf{Open (the Clay Millennium Prize problem):}
\begin{enumerate}
\setcounter{enumi}{4}
\item The continuum-limit hypothesis (H): the limit
$m_{\mathrm{lat}}^{2}/a^{2} \to m_{\mathrm{phys}}^{2} > 0$ as
$a \to 0$, $\beta \to \infty$.
\end{enumerate}

\textbf{Consequently, the Clay Yang--Mills mass gap conjecture is
NOT proved by this repository.} What is proved is the conditional
implication (H) $\Rightarrow$ (Clay mass gap), with the lattice gap
$m_{\mathrm{lat}}^{2} > 0$ rigorously derived from the Wilson action
(replacing the free parameter $\delta = 0.05$). The hypothesis (H)
is a concrete, attackable open problem. See
\texttt{AUDIT\_2026\_07\_06.md} for the full audit and
\texttt{verfification/plaquette\_correlator\_transfer\_operator.py}
for the numerical verification of all the rigorous steps.

\subsection{Comparison with the Ward--Hodge approach}

\begin{center}
\begin{tabular}{l|l|l}
\hline
\textbf{Aspect} & \textbf{Ward--Hodge (app192, app208)} &
\textbf{PCTO (this appendix)}\\
\hline
Lattice gap & Free parameter $\delta = 0.05$ & $m_{\mathrm{lat}}^{2} =
  -2\log\tanh(\beta/N)$ \\
Origin & Chosen & Derived from Wilson action \\
Positivity & Tautological (gap $= \delta$) &
  Rigorous (Theorem \ref{thm:strong_coupling_gap})\\
Weak-coupling behavior & $\rho(K_{W}) \ge 1$ (gap negative) &
  $m_{\mathrm{lat}}^{2} > 0$ for all $\beta > 0$\\
Spectral identification & proved (Hypothesis H.2) &
  Markov intertwining (Theorem \ref{thm:mitt})\\
Continuum transport & proved (Hypothesis H.1) &
  Open (Hypothesis (H) of Theorem \ref{thm:clay_main_H})\\
\hline
\end{tabular}
\end{center}

The PCTO replaces the tautological free parameter with a rigorous
strong-coupling formula, replaces the failing Ward--Hodge spectral
identification with a Markov intertwining theorem, and isolates the
remaining open problem (the continuum-limit hypothesis (H)) as a
concrete, attackable conjecture.

\subsection{Correction (2026-07-07): errors in Theorem~\ref{thm:mitt}
and Theorem~\ref{thm:clay_main_H}}
\label{ssec:pcto_correction}

A further audit identifies two errors in the argument above, corrected
in \S\ref{app:perron_frobenius_mass_gap} (app210):

\begin{enumerate}
\item \textbf{Theorem~\ref{thm:mitt}, step (2), is false as stated.}
The claim ``$T_{\mathrm{OS}}$ preserves $\mathcal{A}_{\mathrm{plaq}}$''
is not justified: the Wilson action couples every plaquette to its
neighbors through shared links, so the conditional expectation
$\Pi=\mathbb{E}[\cdot\mid\mathcal{A}_{\mathrm{plaq}}]$ does not commute
with $T_{\mathrm{OS}}$ except in the trivial case where the $2\times2$
block is the entire lattice (no neighbors to integrate out). The
identity $\Pi T_{\mathrm{OS}}\Pi = T_{\mathrm{OS}}\Pi$ asserted in the
proof requires $T_{\mathrm{OS}}\mathcal{A}_{\mathrm{plaq}}\subseteq
\mathcal{A}_{\mathrm{plaq}}$, which fails once link variables shared
with adjacent plaquettes are integrated out: $T_{\mathrm{OS}}$ maps a
plaquette observable to a function of the boundary links of the whole
time-slice, not just the four plaquettes of the block. The correct
replacement is the Krein--Rutman argument of
\S\ref{app:perron_frobenius_mass_gap}, Theorem
\ref{thm:finite_lattice_mass_gap}, which works directly with
$T_{\mathrm{OS}}$ on the full physical Hilbert space and requires no
projection onto a plaquette sub-algebra.

\item \textbf{Theorem~\ref{thm:clay_main_H}, step (4), is false as
stated.} The claim that
$m_{\mathrm{lat}}^{2}(\beta,N)/a(\beta)^{2}\to m_{\mathrm{phys}}^{2}>0$
under hypothesis (H) is not merely conditional but uses an inapplicable
formula: Lemma~\ref{lem:sc_diverges} of
\S\ref{app:perron_frobenius_mass_gap} shows that the strong-coupling
mass $m^{\mathrm{SC}}(\beta,N)=-2\log\tanh(\beta/N)$, substituted into
the two-loop asymptotic-freedom lattice spacing $a(\beta)$, gives a
ratio $m^{\mathrm{SC}}(\beta,N)^{2}/a(\beta)^{2}\to\infty$ (not a finite
positive limit), because $m^{\mathrm{SC}}$ decays only like
$e^{-4\beta/N}$ while $a(\beta)^{2}$ decays like $e^{-2b_{0}\beta}$
with $2b_{0}=\tfrac{11}{3}N>4/N$ for every $N\ge2$. The strong-coupling
formula is valid only for $\beta<\beta_{c}(N)$
(Theorem~\ref{thm:strong_coupling_gap} is a genuine theorem \emph{in
that regime}); it has no validity and cannot be extrapolated to
$\beta\to\infty$. The corrected hypothesis (H$'$), stated in
\S\ref{app:perron_frobenius_mass_gap}, Theorem~\ref{thm:corrected_H},
replaces $m_{\mathrm{lat}}^{2}(\beta,N)$ by the actual
non-perturbatively defined lattice mass gap $m(\beta)^{2}$ (Definition
\ref{def:actual_mass_gap}), for which no closed-form expression is
known at weak coupling. Theorem~\ref{thm:clay_main_H} should be read as
superseded by Theorem~\ref{thm:corrected_H}.
\end{enumerate}

These corrections do not close the Clay problem; they replace an
erroneous derivation with the honestly open hypothesis (H$'$). See
\S\ref{app:perron_frobenius_mass_gap} for the complete corrected
treatment, including the unconditional finite-lattice mass gap
(Theorem~\ref{thm:finite_lattice_mass_gap}), which is new and does not
appear anywhere else in this repository.

\endinput
